\chapter{Introduction}

Simultaneous Localisation and Mapping (SLAM) is one of the fundamental problems in
mobile robotics~\cite{leonard1991mobile}.
SLAM poses a question: how a robot can build a consistent map of an unknown environment while
simultaneously estimating its own position within that map.
\emph{Localisation} refers to estimating the state of the robot, which
consists of position and orientation of the robot;
\emph{mapping} refers to construction of a model, which is the map of the surrounding environment.
Having a map is useful for two reasons~\cite{cadena2016past}.
First, it enables downstream tasks such as navigation planning and
environment visualisation for a human user.
Second, it keeps localisation errors in check by letting the robot recognise
where it has already been.
SLAM is therefore considered a state of the art for autonomous robots that must
operate without GPS.

The difficulty of the SLAM problem lies in the mutual dependency of its two components.
A good map can only be built if the robot knows its position, and the robot can only
determine its position if it already has a map.
Both must therefore be estimated at the same time from noisy sensor data.
The key mechanism that prevents errors from growing indefinitely is \emph{loop closure}:
when the robot returns to a place it has already visited, it detects the match and uses
it to correct the drift that has built up in its trajectory.
Without it, errors accumulate without bound~\cite{cadena2016past}.

The origins of SLAM research trace back to 1986, when Smith, Self, and Cheeseman~\cite{smith1990estimating},
and independently Durrant-Whyte~\cite{durrantwhyte1988uncertain}, first described geometric uncertainty and the
statistical relationships between robot poses and environmental
landmarks. The term \emph{SLAM} itself was first
used at the International Symposium on Robotics Research (ISRR) in
1995~\cite{durrant2006simultaneous}.
Early work introduced the main
probabilistic formulations: Extended Kalman Filters, Rao-Blackwellised
Particle Filters, and maximum likelihood estimation. 
In following years, research
focused on observability, convergence, and consistency, which produced the
first widely used open-source libraries~\cite{cadena2016past}.
% and established sparsity as the key to efficient solvers. 
Since then, supported sensors
have expanded from 2D laser scanners to 3D LiDARs, cameras, and
tightly-coupled inertial systems, and the current research focus has shifted
towards robust long-term operation.

SLAM has been continuously evolving since its formulation four decades ago, driven largely by the
growing use of mobile robots in indoor environments such as homes or
warehouses, where GPS is not available. Today it is used across many
domains: robotic vacuum cleaners navigate home environments, autonomous
vehicles map urban roads, unmanned aerial vehicles use it for inspection
and surveillance, and underwater drones can monitor coral reefs or detect mines.

% Further applications include underwater reef monitoring, underground mine exploration, and
% terrain mapping for planetary rovers operating beyond GPS
% coverage~\cite{durrant2006simultaneous}.

Despite four decades of progress, SLAM is not a fully solved
problem~\cite{cadena2016past}. It can work reliably for slow-moving robots in
static indoor environments, but can struggle in fast-moving or highly dynamic
outdoor scenarios. A large number of open-source implementations exist, each
taking a different algorithmic approach with its own strengths and
weaknesses. Therefore, selecting the most suitable algorithm for a specific application
requires a direct comparison.

\section{Goal and Scope}

The goal of this thesis is to evaluate and compare SLAM algorithms across different sensor
modalities and algorithmic families under identical, controlled conditions.
The evaluation is carried out on a mobile robot platform equipped with differential drive,
a 3D LiDAR, a stereo camera, and an IMU.
Eleven algorithms are selected to cover the major algorithmic families: sensor fusion,
LiDAR odometry, 2D and 3D LiDAR SLAM, tightly-coupled LiDAR-IMU odometry, and
visual SLAM.

All algorithms are tested on the same dataset recorded from simulated
environment in NVIDIA Isaac Sim, which provides exact ground-truth poses for objective
evaluation.
The framework used for all experiments is ROS\,2 (Robot Operating System), which
provides a common interface for sensor data, algorithm execution, and data recording.

Trajectory accuracy is quantified using Absolute Trajectory Error and Relative Pose Error
against the simulator's ground truth.
Map quality is assessed with Chamfer distance and F-score metrics for algorithms that
produce a 3D point cloud, and visually for those that produce an occupancy grid.
Computational cost is measured as mean CPU usage and peak RAM consumption during
execution.

\section{Thesis Structure}

\begingroup
\setlength{\parskip}{2pt}
\setlength{\parindent}{1.5em}

Chapter~2 contains the theoretical background required to understand the thesis: navigation in mobile
robotics, sensors description, map representations, point cloud registration, and evaluation metrics.

Chapter~3 describes each of the eleven evaluated algorithms, grouped by method family.

Chapter~4 presents the experimental setup, the dataset, and the results for all eleven
algorithms, including trajectory accuracy, map quality, and computational cost.

Chapter~5 concludes the thesis and provides a summary of findings made during experimental
evaluation of algorithms.

\endgroup
